
\chapter{Fresh Ubuntu Setup for Rudra and Aghora}
\label{chap:fresh-setup}

\chapterquote{A robot workspace is only useful when a clean machine can be rebuilt into the same robot again. This chapter is the rebuild contract for the onboard NUC and the workstation.}{RUDRAv4 setup principle}

This chapter turns a fresh Ubuntu installation into a working two-machine \RUDRAv{} development system.  The robot computer is the Intel NUC mounted on the rover and named \cmd{rudra}.  The workstation laptop is named \cmd{aghora}.  Both machines run Ubuntu 26.04 and \ros{} Lyrical, but they have different jobs.  The NUC owns hardware, serial devices, the LiDAR driver, the Teensy bridge, DCDC monitoring, and launch files that move the robot.  Aghora is the engineering workstation: editing, Git, RViz, plotting, bag review, SLAM visualization, and remote operation.

\begin{designbox}{Outcome of the setup}
After this setup, opening a terminal on either machine should show a sourced \ros{} Lyrical environment, a consistent \cmd{ROS_DOMAIN_ID}, and the ability to build the RUDRAv4 workspace.  On \cmd{rudra}, hardware devices should appear under \filep{/dev/serial/by-id}.  On \cmd{aghora}, RViz should discover topics published by the NUC over the robot network.
\end{designbox}

\section{Machine roles}
\begin{table}[h]
\centering
\caption{Two-machine division of responsibility after a fresh install.}
\begin{tabularx}{\linewidth}{p{0.17\linewidth}p{0.28\linewidth}Y}
\toprule
Machine & Identity & Main responsibility \\
\midrule
\cmd{rudra} & Intel NUC7i7BNH on the rover & Hardware-facing \ros{} nodes: Uno bridge, Teensy bridge, LiDAR, obstacle guard, localization, DCDC-USB monitor, logs, and command safety. \\
\cmd{aghora} & Ubuntu workstation/laptop & Development, SSH access, Git, RViz, map building, launch observation, bag analysis, and remote visualization. \\
\bottomrule
\end{tabularx}
\end{table}

Both machines should use the same \cmd{ROS_DOMAIN_ID}.  RUDRAv4 uses \cmd{42} in the current workspace.  A mismatched DDS domain is the most common reason the laptop sees no topics even when both computers are healthy.

\section{Install Ubuntu 26.04}
Install Ubuntu 26.04 on both computers.  The NUC can run Desktop or Server; Desktop is more convenient during development because Arduino IDE, RViz, and VS Code are available locally.  The workstation should use Desktop.

Set the hostnames immediately after the first boot:

\begin{lstlisting}[language=bash,caption={Hostname setup on the NUC.}]
sudo hostnamectl set-hostname rudra
hostnamectl
cat /etc/hostname
\end{lstlisting}

\begin{lstlisting}[language=bash,caption={Hostname setup on the workstation.}]
sudo hostnamectl set-hostname aghora
hostnamectl
cat /etc/hostname
\end{lstlisting}

Update the base operating system on both machines:

\begin{lstlisting}[language=bash,caption={Base Ubuntu update.}]
sudo apt update
sudo apt upgrade -y
sudo apt install -y software-properties-common
sudo add-apt-repository universe
sudo reboot
\end{lstlisting}

\section{Bootstrap packages common to both machines}
The clean install should not assume build tools, Python headers, pip, venv, SSH, or network discovery tools are already present.  Install the common toolchain first:

\begin{lstlisting}[language=bash,caption={Common Linux development packages.}]
sudo apt update
sudo apt install -y \
  build-essential cmake pkg-config ninja-build \
  git curl wget gnupg ca-certificates lsb-release \
  locales software-properties-common \
  python3 python3-dev python3-pip python3-venv python3-setuptools python3-wheel \
  usbutils pciutils net-tools iproute2 udev \
  openssh-server avahi-daemon \
  terminator htop tree tmux unzip zip
\end{lstlisting}

Configure a UTF-8 locale before installing \ros{}:

\begin{lstlisting}[language=bash,caption={UTF-8 locale setup.}]
locale
sudo locale-gen en_US en_US.UTF-8
sudo update-locale LC_ALL=en_US.UTF-8 LANG=en_US.UTF-8
export LANG=en_US.UTF-8
locale
\end{lstlisting}

\section{Install ROS 2 Lyrical}
\ros{} Lyrical is the RUDRAv4 distribution for Ubuntu 26.04.  The official Lyrical installation page lists Ubuntu Resolute 26.04 as the binary package target.  Use the official ROS apt-source package rather than hand-written apt source lines because it avoids the signed-by conflicts that occur when old ROS source entries are left behind.

\begin{lstlisting}[language=bash,caption={Install the ROS apt source.}]
sudo apt update
sudo apt install -y curl
export ROS_APT_SOURCE_VERSION=$(curl -s https://api.github.com/repos/ros-infrastructure/ros-apt-source/releases/latest | grep -F "tag_name" | awk -F'"' '{print $4}')
curl -L -o /tmp/ros2-apt-source.deb \
  "https://github.com/ros-infrastructure/ros-apt-source/releases/download/${ROS_APT_SOURCE_VERSION}/ros2-apt-source_${ROS_APT_SOURCE_VERSION}.$(. /etc/os-release && echo ${UBUNTU_CODENAME:-${VERSION_CODENAME}})_all.deb"
sudo dpkg -i /tmp/ros2-apt-source.deb
\end{lstlisting}

Install the desktop distribution on both machines during development.  If the NUC must later be reduced, it can use \cmd{ros-lyrical-ros-base}, but Desktop is easier while bringing up a robot.

\begin{lstlisting}[language=bash,caption={Install ROS 2 Lyrical and core tools.}]
sudo apt update
sudo apt upgrade -y
sudo apt install -y \
  ros-lyrical-desktop \
  ros-dev-tools \
  python3-colcon-common-extensions \
  python3-rosdep \
  python3-vcstool
\end{lstlisting}

Install packages used by the current RUDRAv4 stack:

\begin{lstlisting}[language=bash,caption={RUDRAv4 ROS package set.}]
sudo apt install -y \
  ros-lyrical-rmw-fastrtps-cpp \
  ros-lyrical-rmw-cyclonedds-cpp \
  ros-lyrical-robot-state-publisher \
  ros-lyrical-joint-state-publisher-gui \
  ros-lyrical-xacro \
  ros-lyrical-tf2-tools \
  ros-lyrical-teleop-twist-keyboard \
  ros-lyrical-teleop-twist-joy \
  ros-lyrical-joy \
  ros-lyrical-joy-linux \
  ros-lyrical-imu-tools \
  ros-lyrical-imu-filter-madgwick \
  ros-lyrical-imu-complementary-filter \
  ros-lyrical-rviz-imu-plugin \
  ros-lyrical-robot-localization \
  diagnostic-updater \
  python3-serial \
  libcxx-serial-dev \
  libserialport-dev
\end{lstlisting}

Initialize rosdep:

\begin{lstlisting}[language=bash,caption={rosdep initialization.}]
sudo rosdep init || true
rosdep update
\end{lstlisting}

Verify ROS itself before cloning RUDRA:

\begin{lstlisting}[language=bash,caption={Basic ROS verification.}]
source /opt/ros/lyrical/setup.bash
echo "ROS_DISTRO=$ROS_DISTRO"
which ros2
ros2 doctor --report
\end{lstlisting}

\section{Clone and build the workspace}
On the NUC, the current workspace is:

\begin{lstlisting}[language=bash,caption={RUDRAv4 workspace path on the NUC.}]
mkdir -p /home/rudra/Projects
cd /home/rudra/Projects
git clone https://github.com/AdarshGouda/RUDRAv4.git
cd /home/rudra/Projects/RUDRAv4
\end{lstlisting}

On Aghora, use the same repository path when possible.  It keeps scripts, documentation, and local notes consistent:

\begin{lstlisting}[language=bash,caption={RUDRAv4 workspace path on Aghora.}]
mkdir -p ~/Projects
cd ~/Projects
git clone https://github.com/AdarshGouda/RUDRAv4.git
cd ~/Projects/RUDRAv4
\end{lstlisting}

Build the RUDRAv4 workspace:

\begin{lstlisting}[language=bash,caption={Build RUDRAv4.}]
cd /home/rudra/Projects/RUDRAv4
source /opt/ros/lyrical/setup.bash
rosdep install --from-paths src --ignore-src -r -y
colcon build --symlink-install
source install/setup.bash
ros2 pkg list | grep rudra
\end{lstlisting}

\section{Install serial permissions and board tools on the NUC}
The NUC opens three important USB serial devices: Arduino Uno, Teensy 4.1, and the CP2102 USB adapter used by the YDLIDAR.  Add the user to the serial groups and install udev rules before testing the robot.

\begin{lstlisting}[language=bash,caption={Serial permissions on the NUC.}]
sudo usermod -aG dialout,plugdev rudra
groups rudra
# Log out and back in after this step.
\end{lstlisting}

Install Teensy rules:

\begin{lstlisting}[language=bash,caption={Teensy udev rules.}]
curl -L -o /tmp/00-teensy.rules https://www.pjrc.com/teensy/00-teensy.rules
sudo install -m 0644 /tmp/00-teensy.rules /etc/udev/rules.d/00-teensy.rules
sudo udevadm control --reload-rules
sudo udevadm trigger
\end{lstlisting}

Install Arduino CLI and Teensy board support:

\begin{lstlisting}[language=bash,caption={Arduino CLI and Teensy board index.}]
sudo apt install -y arduino-cli
arduino-cli config init || true
arduino-cli config add board_manager.additional_urls https://www.pjrc.com/teensy/package_teensy_index.json
arduino-cli core update-index
arduino-cli core search teensy
\end{lstlisting}

Validate serial devices:

\begin{lstlisting}[language=bash,caption={Serial device identification.}]
cd /home/rudra/Projects/RUDRAv4
bash scripts/list_serial.sh
lsusb
ls -l /dev/serial/by-id
\end{lstlisting}

\section{Install the YDLIDAR external workspace on the NUC}
The current RUDRAv4 workspace treats the LiDAR driver as an external ROS workspace under \filep{/home/rudra/ros2_ws}.  Source order matters: base ROS, external LiDAR workspace, then RUDRAv4.

\begin{lstlisting}[language=bash,caption={Create the external LiDAR workspace.}]
mkdir -p /home/rudra/ros2_ws/src
cd /home/rudra/ros2_ws/src
git clone https://github.com/YDLIDAR/YDLidar-SDK.git
git clone -b humble https://github.com/YDLIDAR/ydlidar_ros2_driver.git
\end{lstlisting}

Build the SDK into the local external workspace rather than installing it globally:

\begin{lstlisting}[language=bash,caption={Build YDLidar-SDK locally.}]
cd /home/rudra/ros2_ws/src/YDLidar-SDK
mkdir -p build
cd build
cmake .. -DCMAKE_INSTALL_PREFIX=/home/rudra/ros2_ws/install/YDLidar-SDK
cmake --build .
cmake --install .
\end{lstlisting}

Build the ROS driver:

\begin{lstlisting}[language=bash,caption={Build YDLIDAR ROS 2 driver against the local SDK.}]
cd /home/rudra/ros2_ws
source /opt/ros/lyrical/setup.bash
export CMAKE_PREFIX_PATH=/home/rudra/ros2_ws/install/YDLidar-SDK:$CMAKE_PREFIX_PATH
colcon build --symlink-install --packages-select ydlidar_ros2_driver
source install/setup.bash
\end{lstlisting}

\section{Install the voice stack on the NUC}
\pkg{rudra_voice} needs a small Python environment that is separate from system Python but still able to see ROS Python packages. This is the fix for the \cmd{No module named 'sounddevice'} failure seen when ROS launches the node with the wrong interpreter.

Install audio and Python system dependencies:

\begin{lstlisting}[language=bash,caption={Voice OS dependencies.}]
sudo apt update
sudo apt install -y \
  alsa-utils espeak-ng \
  python3-venv python3-pip \
  libportaudio2 portaudio19-dev \
  curl wget unzip
\end{lstlisting}

Create the project voice virtual environment:

\begin{lstlisting}[language=bash,caption={Voice Python virtual environment.}]
cd /home/rudra/Projects/RUDRAv4
python3 -m venv --system-site-packages .venv_voice
source .venv_voice/bin/activate
python -m pip install --upgrade pip
python -m pip install vosk sounddevice numpy requests
python -c "import sounddevice, vosk, requests; print('voice python ok')"
\end{lstlisting}

\begin{designbox}{Why \cmd{--system-site-packages} is intentional}
ROS 2 Lyrical installs important Python modules through apt. The voice stack installs speech packages through pip. The \filep{.venv_voice} environment with \cmd{--system-site-packages} lets one Python interpreter see both worlds. If pip reports \cmd{externally-managed-environment}, pip is pointed at the system interpreter instead of \filep{.venv_voice}.
\end{designbox}

Download the Vosk speech-recognition model:

\begin{lstlisting}[language=bash,caption={Vosk model install.}]
mkdir -p /home/rudra/Projects/RUDRAv4/models
cd /home/rudra/Projects/RUDRAv4/models
wget https://alphacephei.com/vosk/models/vosk-model-small-en-us-0.15.zip
unzip vosk-model-small-en-us-0.15.zip
\end{lstlisting}

Install Piper locally for human-sounding offline speech. This keeps the install reproducible even when \cmd{sudo apt install piper} is not available or requires an interactive password:

\begin{lstlisting}[language=bash,caption={Local Piper install.}]
cd /home/rudra/Projects/RUDRAv4
mkdir -p tools/piper models/piper
curl -L --fail -o /tmp/piper_linux_x86_64.tar.gz \
  https://github.com/rhasspy/piper/releases/download/2023.11.14-2/piper_linux_x86_64.tar.gz
tar -xzf /tmp/piper_linux_x86_64.tar.gz -C tools/piper --strip-components=1
curl -L --fail -o models/piper/en_US-amy-medium.onnx \
  https://huggingface.co/rhasspy/piper-voices/resolve/main/en/en_US/amy/medium/en_US-amy-medium.onnx
curl -L --fail -o models/piper/en_US-amy-medium.onnx.json \
  https://huggingface.co/rhasspy/piper-voices/resolve/main/en/en_US/amy/medium/en_US-amy-medium.onnx.json
printf 'I am online. Ready for safe commands.\n' | \
  tools/piper/piper --model models/piper/en_US-amy-medium.onnx \
  --output_file /tmp/rudra_piper_test.wav
aplay /tmp/rudra_piper_test.wav
\end{lstlisting}

Check the P610 microphone and speakerphone:

\begin{lstlisting}[language=bash,caption={P610 audio check.}]
lsusb
arecord -l
aplay -l
arecord -f S16_LE -r 16000 -c 1 -d 5 /tmp/rudra_p610_test.wav
aplay /tmp/rudra_p610_test.wav
\end{lstlisting}

Ollama is optional. Install it only if natural-language routing is needed beyond deterministic commands. The full robot launch can start or check Ollama for you when \cmd{enable_ollama:=true}; the helper script below is also useful as a one-time readiness check after installation.

\begin{lstlisting}[language=bash,caption={Optional Ollama setup.}]
curl -fsSL https://ollama.com/install.sh | sh
bash /home/rudra/Projects/RUDRAv4/scripts/start_ollama.sh qwen2.5:3b
\end{lstlisting}

\section{Shell environment for every terminal}
Add the following block to \filep{~/.bashrc} on \cmd{rudra}.  On \cmd{aghora}, omit the external LiDAR workspace if it is not cloned locally.

\begin{lstlisting}[language=bash,caption={Recommended NUC bash environment.}]
# RUDRA ROS 2 environment
export RUDRA_WS="$HOME/Projects/RUDRAv4"
export RUDRA_LIDAR_WS="$HOME/ros2_ws"

if [ -f /opt/ros/lyrical/setup.bash ]; then
    source /opt/ros/lyrical/setup.bash
fi
if [ -f "$RUDRA_LIDAR_WS/install/setup.bash" ]; then
    source "$RUDRA_LIDAR_WS/install/setup.bash"
fi
if [ -f "$RUDRA_WS/install/setup.bash" ]; then
    source "$RUDRA_WS/install/setup.bash"
fi

export ROS_DOMAIN_ID=42
export ROS_LOCALHOST_ONLY=0
export ROS_AUTOMATIC_DISCOVERY_RANGE=SUBNET
export RMW_IMPLEMENTATION=rmw_fastrtps_cpp

alias rudra_ws='cd $RUDRA_WS'
alias cb='cd $RUDRA_WS && colcon build --symlink-install'
alias serialids='ls -l /dev/serial/by-id 2>/dev/null'
alias rosenv="printenv | grep -E 'ROS|RMW|RUDRA'"
printf '\n[RUDRA] ROS 2 Lyrical environment loaded. Workspace: %s\n' "$RUDRA_WS"
\end{lstlisting}

\section{SSH and network verification}
Enable SSH and mDNS on the NUC:

\begin{lstlisting}[language=bash,caption={NUC remote access.}]
sudo systemctl enable --now ssh
sudo systemctl enable --now avahi-daemon
hostname -I
systemctl status ssh --no-pager
\end{lstlisting}

From Aghora:

\begin{lstlisting}[language=bash,caption={SSH from Aghora to Rudra.}]
ssh rudra@rudra.local
\end{lstlisting}

Optional \filep{~/.ssh/config} entry on Aghora:

\begin{lstlisting}[caption={Aghora SSH shortcut.}]
Host rudra
    HostName rudra.local
    User rudra
    IdentityFile ~/.ssh/id_ed25519
\end{lstlisting}

Test ROS discovery.  On \cmd{rudra}:

\begin{lstlisting}[language=bash,caption={Publish a topic from the NUC.}]
source ~/.bashrc
ros2 run demo_nodes_cpp talker
\end{lstlisting}

On \cmd{aghora}:

\begin{lstlisting}[language=bash,caption={Listen from the workstation.}]
source ~/.bashrc
ros2 topic list
ros2 topic echo /chatter
\end{lstlisting}

\section{First clean-room bring-up checklist}
\begin{enumerate}
  \item Ubuntu updated on both machines.
  \item \cmd{rudra} and \cmd{aghora} hostnames set.
  \item \ros{} Lyrical installed and sourced on both machines.
  \item \cmd{ROS_DOMAIN_ID=42} on both machines.
  \item RUDRAv4 workspace cloned and built.
  \item \filep{.venv_voice} exists and \cmd{python -c "import sounddevice, vosk, requests"} succeeds inside it.
  \item Vosk model exists under \filep{models/vosk-model-small-en-us-0.15}.
  \item Piper binary exists at \filep{tools/piper/piper}; Piper test audio plays through \cmd{aplay}.
  \item \cmd{espeak-ng} is installed as the fallback speech output.
  \item NUC user belongs to \cmd{dialout} and \cmd{plugdev}.
  \item Teensy rules installed.
  \item \filep{/dev/serial/by-id} shows Uno, Teensy, and LiDAR adapter.
  \item Aghora can SSH into \cmd{rudra}.
  \item Aghora can see ROS topics from \cmd{rudra}.
\end{enumerate}

\section{First integrated launch after clean setup}
After the checklist passes, use the single full launch file:

\begin{lstlisting}[language=bash,caption={Integrated RUDRAv4 bringup on the NUC.}]
cd /home/rudra/Projects/RUDRAv4
source /opt/ros/lyrical/setup.bash
source /home/rudra/ros2_ws/install/setup.bash
source install/setup.bash
ros2 launch rudra_base_bridge rudra_full.launch.py
\end{lstlisting}

For the normal robot-mounted NUC, enable DCDC monitoring:

\begin{lstlisting}[language=bash,caption={Integrated launch with NUC power monitoring.}]
ros2 launch rudra_base_bridge rudra_full.launch.py enable_dcdc_monitor:=true
\end{lstlisting}

To let the voice stack use Ollama for uncertain phrases, opt in explicitly:

\begin{lstlisting}[language=bash,caption={Integrated launch with Ollama voice routing.}]
ros2 launch rudra_base_bridge rudra_full.launch.py \
  enable_dcdc_monitor:=true \
  enable_ollama:=true \
  voice_use_llm_router:=true
\end{lstlisting}

Expected early launch signs are: the launch prints the voice Python as \filep{.venv_voice/bin/python}, the P610 is detected by name, \topic{/scan} appears after the LiDAR finishes startup, \topic{/rudra/wheel_odom} and \topic{/rudra/imu/raw} publish from Teensy telemetry, \topic{/odometry/filtered} publishes from the EKF, and \topic{/rudra_voice/transcript} updates when speech is heard. The full launch keeps \cmd{voice_use_llm_router} off unless you turn it on, so deterministic voice commands work even if Ollama is stopped.

\begin{operatorbox}{First launch rule of thumb}
After a clean NUC install, use \cmd{ros2 launch rudra_base_bridge rudra_full.launch.py} as the main bringup command. Add launch arguments only for the capability being tested: \cmd{enable_dcdc_monitor:=true} for NUC power diagnostics, \cmd{enable_rviz:=true} for local visualization, and \cmd{enable_ollama:=true voice_use_llm_router:=true} for LLM-assisted voice phrases.
\end{operatorbox}

\begin{safetybox}{Do not test motion yet}
The purpose of this chapter is to rebuild the software platform.  Wheels should stay lifted or motor battery power should stay disconnected until the firmware and serial bridge chapters have been validated.
\end{safetybox}
